\documentclass[../../main.tex]{subfiles}% 注意这里的文件路径不能用 ./main.tex ,否则用latexmk编译子文件会报错
\graphicspath{{\subfix{./image/}}} % 指定图片目录,后续可以直接使用图片文件名
% 注意这里的文件路径不能用 ../../image/ ,否则用latexmk编译子文件会报错

% 例如:
% \begin{figure}[H]
% \centering
% \includegraphics[scale=0.3]{图.png}
% \caption{}
% \label{figure:图}
% \end{figure}
% 注意:上述\label{}一定要放在\caption{}之后,否则引用图片序号会只会显示??.

\begin{document}

\section{有限域}

\begin{theorem}
对任意素数的幂$p^n$，在同构意义下存在唯一的有限域$F$，其元素个数为$p^n$。
\end{theorem}
\begin{proof}
设有限域$F$的元素个数为$p^n$，则其非零元素集合$F^*=F\setminus\{0\}$构成阶为$p^n-1$的乘法群。
对任意$\alpha\in F^*$，由拉格朗日定理得$\alpha^{p^n-1}=1$，故对任意$\alpha\in F$，均为$\mathbb{Z}_p[x]$中多项式$f(x)=x^{p^n}-x$的根。
求导得$f'(x)=-1$，故$f(x)$无重根，恰有$p^n$个根。

设$f(\alpha)=f(\beta)=0$，则
\begin{align*}
(\alpha+\beta)^{p^n}=\alpha^{p^n}+\beta^{p^n}=\alpha+\beta,\\
(\alpha\beta)^{p^n}=\alpha^{p^n}\beta^{p^n}=\alpha\beta,\\
(\alpha^{-1})^{p^n}=(\alpha^{p^n})^{-1}=\alpha^{-1},\quad \alpha\neq0.
\end{align*}
故$f(x)$的根集构成$p^n$元域，且为$f(x)$的分裂域。由分裂域的同构唯一性，$F$在同构意义下唯一。

\end{proof}

\begin{remark}
由有限交换群性质，$F^*$为循环群，其生成元$\xi$是$p^n-1$次本原单位根，此时$F=\mathbb{Z}_p(\xi)$，将该有限域记为$F_{p^n}$。
\end{remark}

\begin{theorem}
设$p$是素数，$n$是自然数，则$\mathrm{Gal}(F_{p^n}/\mathbb{Z}_p)$是$n$阶循环群。
\end{theorem}
\begin{proof}
定义映射$\sigma:F_{p^n}\to F_{p^n},\ \sigma(\alpha)=\alpha^p$，易证$\sigma\in\mathrm{Gal}(F_{p^n}/\mathbb{Z}_p)$。
由$F_{p^n}$中元素满足$\alpha^{p^n}=\alpha$，得$\sigma^n=\mathrm{id}$，即$\sigma$生成一个$n$阶循环群。

因$F_{p^n}=\mathbb{Z}_p(\xi)=\{0,\xi^i\mid 1\leqslant i\leqslant p^n-1\}$，考虑多项式
\begin{align*}
f(x)=\prod\limits_{i=0}^{n-1}(x-\xi^{p^i}).
\end{align*}
$f(x)$的系数为$\xi,\xi^p,\dots,\xi^{p^{n-1}}$的对称多项式，任取系数$g=g(\xi,\xi^p,\dots,\xi^{p^{n-1}})$，有
\begin{align*}
g^p=g(\xi^p,\dots,\xi^{p^{n-1}},\xi)=g(\xi,\xi^p,\dots,\xi^{p^{n-1}})=g.
\end{align*}
因$x^p=x$在$\mathbb{Z}_p$中仅有$p$个根，故$g\in\mathbb{Z}_p$，即$f(x)\in\mathbb{Z}_p[x]$。

对任意$\theta\in\mathrm{Gal}(F_{p^n}/\mathbb{Z}_p)$，有
\begin{align*}
\theta(f(x))=f(x)=\prod\limits_{i=0}^{n-1}(x-\theta(\xi)^{p^i}),
\end{align*}
故$\theta(\xi)=\xi^{p^k}$，即$\theta=\sigma^k$。因此$\mathrm{Gal}(F_{p^n}/\mathbb{Z}_p)=\langle\sigma\rangle$为$n$阶循环群。

\end{proof}

\begin{example}
设$p$是奇素数，则
\begin{enumerate}[(1)]
\item $\exists\alpha\in\mathbb{Z}_p$，使得$\alpha^2=-1$当且仅当$p\equiv1\pmod{4}$；
\item $p$可表为两个整数的平方和当且仅当$p\equiv1\pmod{4}$。
\end{enumerate}
\end{example}
\begin{proof}
(1) 循环群$\mathbb{Z}_p^*=\langle\xi\rangle$的阶为$p-1$。令$\beta=\xi^{\frac{p-1}{2}}$，则$\beta\neq1$且$\beta^2=1$，故$\beta=-1$。
若$p\equiv1\pmod{4}$，取$\alpha=\xi^{\frac{p-1}{4}}$，则$\alpha^2=-1$。
若$p\equiv3\pmod{4}$，则$p-1=2m$（$m$为奇数），对任意$k$，$4\nmid\mathrm{ord}(\xi^k)$，故$(\xi^k)^2\neq-1$。

(2) 若$p\equiv3\pmod{4}$，假设$p=m^2+n^2\ (0<m,n<p)$，则在$\mathbb{Z}_p$中$\overline{m}^2+\overline{n}^2=\overline{0}$，即$(\overline{m}\overline{n}^{-1})^2=-1$，与(1)矛盾。
若$p\equiv1\pmod{4}$，由(1)，存在$a\in\mathbb{Z}$使得在$\mathbb{Z}_p$中$\overline{a}^2=-1$，即$p\mid a^2+1$。
因$R=\mathbb{Z}[\sqrt{-1}]$是主理想整环，理想$I=\langle p,a+\sqrt{-1}\rangle=\langle\gamma\rangle$，其中$\gamma=mp+n(a+\sqrt{-1})$。
计算得
\begin{align*}
|\gamma|^2=(mp+na)^2+n^2=m^2p^2+n^2(a^2+1)^2+2mnpa,
\end{align*}
故$p\mid|\gamma|^2$。对任意$\delta\in I$，$p\mid|\delta|^2$，故$1\notin I$。
又$a+\sqrt{-1}\notin\langle p\rangle$，故$I\neq R,\langle p\rangle$。
因$p\in I$，设$p=\gamma\gamma_1$，则$|\gamma_1|>1$。由$p^2=|\gamma|^2|\gamma_1|^2$，得$|\gamma|^2=|\gamma_1|^2=p$，故$p=(mp+na)^2+n^2$。

\end{proof}

\begin{definition}
设$K$是体，$V$是加法交换群，存在映射$K\times V\to V,\ (\alpha,v)\mapsto\alpha v$满足
\begin{enumerate}[(1)]
\item $\alpha(v_1+v_2)=\alpha v_1+\alpha v_2$；
\item $(\alpha+\beta)v=\alpha v+\beta v$；
\item $(\alpha\beta)v=\alpha(\beta v)$；
\item $1v=v$，
\end{enumerate}
则称$V$是$K$上的\textbf{左线性空间}。

同理可定义$K$上的\textbf{右线性空间}$U$，映射$U\times K\to U,\ (u,\alpha)\mapsto u\alpha$满足类似四条公理。
\end{definition}
\begin{remark}
体上左、右线性空间具有与域上线性空间一致的线性相关、线性无关、秩、维数、基、子空间、线性映射、线性变换等概念与结论。
$K$上左(右)线性空间$V$中，向量组$v_1,v_2,\dots,v_s$生成的子空间记为$L(v_1,v_2,\dots,v_s)$。
\end{remark}

\begin{example}
在$K^{1\times n}=\{(\alpha_1,\alpha_2,\dots,\alpha_n)\mid\alpha_i\in K\}$中定义加法与数乘：
\begin{align*}
(\alpha_1,\alpha_2,\dots,\alpha_n)+(\beta_1,\beta_2,\dots,\beta_n)=(\alpha_1+\beta_1,\alpha_2+\beta_2,\dots,\alpha_n+\beta_n),\\
\alpha(\alpha_1,\alpha_2,\dots,\alpha_n)=(\alpha\alpha_1,\alpha\alpha_2,\dots,\alpha\alpha_n),
\end{align*}
则$K^{1\times n}$是$K$上的$n$维左线性空间，称为\textbf{$n$维行向量空间}。
$K$上任意$n$维左线性空间均与$K^{1\times n}$同构。
\end{example}

\begin{example}
在$K^{n\times1}=\left\{\begin{pmatrix}\alpha_1\\\alpha_2\\\vdots\\\alpha_n\end{pmatrix}\,\bigg|\,\alpha_i\in K\right\}$中定义加法与数乘：
\begin{align*}
\begin{pmatrix}\alpha_1\\\alpha_2\\\vdots\\\alpha_n\end{pmatrix}+\begin{pmatrix}\beta_1\\\beta_2\\\vdots\\\beta_n\end{pmatrix}=\begin{pmatrix}\alpha_1+\beta_1\\\alpha_2+\beta_2\\\vdots\\\alpha_n+\beta_n\end{pmatrix},\\
\begin{pmatrix}\alpha_1\\\alpha_2\\\vdots\\\alpha_n\end{pmatrix}\alpha=\begin{pmatrix}\alpha_1\alpha\\\alpha_2\alpha\\\vdots\\\alpha_n\alpha\end{pmatrix},
\end{align*}
则$K^{n\times1}$是$K$上的$n$维右线性空间，称为\textbf{$n$维列向量空间}。
$K$上任意$n$维右线性空间均与$K^{n\times1}$同构。
\end{example}

\begin{example}
设$L$是体$K$的子体，则$K$是$L$上的左(右)线性空间，其维数记为$[K:L]$。
特别地，若$K$是有限体，则$[K:L]<\infty$。
\end{example}

\begin{theorem}
设体满足$M\subseteq L\subseteq K$，则
\begin{align*}
[K:M]=[K:L][L:M].
\end{align*}
\end{theorem}
\begin{proof}
仅证明有限维情形。
设$v_1,v_2,\dots,v_m$是$K$对$L$的基，即$m=[K:L]$；$u_1,u_2,\dots,u_n$是$L$对$M$的基，即$n=[L:M]$。

对任意$\alpha\in K$，有
\begin{align*}
\alpha=\sum\limits_{k=1}^{m}a_kv_k,\quad a_k\in L.
\end{align*}
对每个$a_k\in L$，有
\begin{align*}
a_k=\sum\limits_{j=1}^{n}a_{kj}u_j,\quad a_{kj}\in M.
\end{align*}
代入得
\begin{align*}
\alpha=\sum\limits_{k=1}^{m}\sum\limits_{j=1}^{n}a_{kj}u_jv_k.
\end{align*}

若
\begin{align*}
\sum\limits_{k=1}^{m}\sum\limits_{j=1}^{n}a_{kj}u_jv_k=0,
\end{align*}
由$v_1,\dots,v_m$是$K$对$L$的基，得
\begin{align*}
\sum\limits_{j=1}^{n}a_{kj}u_j=0,\quad 1\leqslant k\leqslant m.
\end{align*}
再由$u_1,\dots,u_n$是$L$对$M$的基，得
\begin{align*}
a_{kj}=0,\quad 1\leqslant k\leqslant m,\ 1\leqslant j\leqslant n.
\end{align*}
故$\{u_jv_k\mid 1\leqslant k\leqslant m,1\leqslant j\leqslant n\}$是$K$对$M$的基，因此$[K:M]=mn=[K:L][L:M]$。

\end{proof}

\begin{theorem}
(Wedderburn) 有限体必为域。
\end{theorem}
\begin{proof}
设$Z$是有限体$K$的中心，即$Z=\{x\in K\mid xy=yx,\forall y\in K\}$，则$Z$是域。记$|Z|=q$，$K$是$Z$上$n$维线性空间，故$|K|=q^n$。
$K^*=K\setminus\{0\}$是$q^n-1$阶群。

任取$a\in K^*$，令$K_a=\{x\in K\mid ax=xa\}$，则$K_a$是$K$的子体且$K_a\supseteq Z$。设$[K_a:Z]=d$，由维数塔定理得$d\mid n$，故$|K_a|=q^d$。
$K_a^*=K_a\setminus\{0\}$是$a$在$K^*$中的中心化子，故$a$的共轭类$C_a$的元素个数为$\frac{q^n-1}{q^d-1}$。

若$n>1$，由群的共轭类分解
\begin{align*}
K^*=Z^*\cup C_{a_1}\cup\cdots\cup C_{a_k},\quad Z^*=Z\setminus\{0\},
\end{align*}
得
\begin{align*}
q^n-1=q-1+\sum\limits_{i=1}^{k}\frac{q^n-1}{q^{d_i}-1},\quad d_i\mid n,\ |C_{a_i}|=\frac{q^n-1}{q^{d_i}-1}.
\end{align*}

由分圆多项式性质，$\Phi_n(x)\,\bigg|\,\frac{x^n-1}{x^{d_i}-1}$，故$\Phi_n(q)\,\bigg|\,\frac{q^n-1}{q^{d_i}-1}$，进而$\Phi_n(q)\mid q-1$。

设$\theta_j\ (1\leqslant j\leqslant\varphi(n))$为$n$次本原单位根，则
\begin{align*}
|\Phi_n(q)|=\prod\limits_{j=1}^{\varphi(n)}|q-\theta_j|>(q-1)^{\varphi(n)}\geqslant q-1,
\end{align*}
矛盾。故$n=1$，即$K=Z$，$K$为域。

\end{proof}



\end{document}